\subsection{2.4  AI Governance and Documentation Requirements}

Regulatory frameworks for AI systems are increasingly emphasizing documentation, human oversight, and traceability. The EU AI Act [4] requires providers of high-risk AI systems to implement post-market monitoring systems and maintain logs sufficient to identify sources of error. US executive orders on AI safety have similarly emphasized the importance of human oversight in consequential automated decision systems. Brundage et al. [3] argue that verifiable claims require mechanisms beyond pre-deployment testing alone. At FAccT 2024, Kolt, Heim, and Anderljung [26] identified three visibility mechanisms for AI agents---identifiers, real-time monitoring, and activity logging---arguing accountability requires linking actions to specific agents and operators. The APIP schema's apip\_id, agent\_id, owner, and exit\_evaluated\_at fields operationalize this linkage, positioning the APIP as both a technical practice and a governance instrument.

